% GENERATED FILE. Edit canonical lessons, then run npm run generate:books.
% canonical-source-hash: fnv1a64:af35dbafdde3e431
% canonical-lessons: BN-C154-rog, BN-C154-sordi, BN-C154-kashi, BN-C154-mathabtha, BN-C154-petbtha

\chapter{Disease, Cold, Cough, Headache, Stomach ache}
\label{ch:bn-disease-cold-cough-headache-stomach-ache}
\hlchaptermodality{154}

\begin{chapteropening}
\textbf{By the end of this chapter:} \emph{I can say a disease, a cold, a cough, a headache and a stomach ache.}

Complete the last lesson of chapter 154: I can say a disease, a cold, a cough, a headache and a stomach ache.
\end{chapteropening}

\section[\texorpdfstring{\bn{রোগ} (rog)}{রোগ (rog)}]{\bn{রোগ} (rog) — a disease}

\label{lesson:BN-C154-rog}



\begin{quote}
Before the new one: say the Bengali for to thank, then the Bengali for to put up with.
\end{quote}

\subsection*{You'll want to know: \bn{রোগ}}
\textbf{\bn{রোগ}} — \emph{rog} — \textquotedblleft{}a disease\textquotedblright{}.

\begin{grammarlens}[title={What you've built}]
One more word for this chapter.
\end{grammarlens}

\subsection*{Your turn}
\begin{itemize}
\raggedright
  \item \emph{Say it:} \emph{rog}
  \item \emph{Say it:} \emph{rog}, once more
  \item \emph{Say it:} say \emph{rog}
  \item \emph{Recall:} say the Bengali for to thank, then the Bengali for to put up with, then say \emph{rog} again
\end{itemize}

\begin{tcolorbox}[breakable,colback=teal!4,colframe=teal!35!black,title={Before you move on}]
What does \bn{রোগ} mean? (\textquotedblleft{}A disease\textquotedblright{}.) Say it once more.
\end{tcolorbox}
\section[\texorpdfstring{\bn{সর্দি} (sôrdi)}{সর্দি (sôrdi)}]{\bn{সর্দি} (sôrdi) — a cold (illness)}

\label{lesson:BN-C154-sordi}



\begin{quote}
Before the new one: say the Bengali for to put up with, then the Bengali for a disease.
\end{quote}

\subsection*{You'll want to know: \bn{সর্দি}}
\textbf{\bn{সর্দি}} — \emph{sôrdi} — \textquotedblleft{}a cold (illness)\textquotedblright{}.

\begin{grammarlens}[title={What you've built}]
One more word for this chapter.
\end{grammarlens}

\subsection*{Your turn}
\begin{itemize}
\raggedright
  \item \emph{Say it:} \emph{sôrdi}
  \item \emph{Say it:} \emph{sôrdi}, once more
  \item \emph{Say it:} say \emph{sôrdi}
  \item \emph{Recall:} say the Bengali for to put up with, then the Bengali for a disease, then say \emph{sôrdi} again
\end{itemize}

\begin{tcolorbox}[breakable,colback=teal!4,colframe=teal!35!black,title={Before you move on}]
What does \bn{সর্দি} mean? (\textquotedblleft{}A cold (illness)\textquotedblright{}.) Say it once more.
\end{tcolorbox}
\section[\texorpdfstring{\bn{কাশি} (kāshi)}{কাশি (kāshi)}]{\bn{কাশি} (kāshi) — a cough}

\label{lesson:BN-C154-kashi}



\begin{quote}
Before the new one: say the Bengali for a disease, then the Bengali for a cold.
\end{quote}

\subsection*{You'll want to know: \bn{কাশি}}
\textbf{\bn{কাশি}} — \emph{kāshi} — \textquotedblleft{}a cough\textquotedblright{}.

\begin{grammarlens}[title={What you've built}]
One more word for this chapter.
\end{grammarlens}

\subsection*{Your turn}
\begin{itemize}
\raggedright
  \item \emph{Say it:} \emph{kāshi}
  \item \emph{Say it:} \emph{kāshi}, once more
  \item \emph{Say it:} say \emph{kāshi}
  \item \emph{Recall:} say the Bengali for a disease, then the Bengali for a cold, then say \emph{kāshi} again
\end{itemize}

\begin{tcolorbox}[breakable,colback=teal!4,colframe=teal!35!black,title={Before you move on}]
What does \bn{কাশি} mean? (\textquotedblleft{}A cough\textquotedblright{}.) Say it once more.
\end{tcolorbox}
\section[\texorpdfstring{\bn{মাথাব্যথা} (māthābæthā)}{মাথাব্যথা (māthābæthā)}]{\bn{মাথাব্যথা} (māthābæthā) — a headache}

\label{lesson:BN-C154-mathabtha}



\begin{quote}
Before the new one: say the Bengali for a cold, then the Bengali for a cough.
\end{quote}

\subsection*{You'll want to know: \bn{মাথাব্যথা}}
\textbf{\bn{মাথাব্যথা}} — \emph{māthābæthā} — \textquotedblleft{}a headache\textquotedblright{}.

\begin{grammarlens}[title={What you've built}]
One more word for this chapter.
\end{grammarlens}

\subsection*{Your turn}
\begin{itemize}
\raggedright
  \item \emph{Say it:} \emph{māthābæthā}
  \item \emph{Say it:} \emph{māthābæthā}, once more
  \item \emph{Say it:} say \emph{māthābæthā}
  \item \emph{Recall:} say the Bengali for a cold, then the Bengali for a cough, then say \emph{māthābæthā} again
\end{itemize}

\begin{tcolorbox}[breakable,colback=teal!4,colframe=teal!35!black,title={Before you move on}]
What does \bn{মাথাব্যথা} mean? (\textquotedblleft{}A headache\textquotedblright{}.) Say it once more.
\end{tcolorbox}
\section[\texorpdfstring{\bn{পেটব্যথা} (peṭbæthā)}{পেটব্যথা (peṭbæthā)}]{\bn{পেটব্যথা} (peṭbæthā) — a stomach ache}

\label{lesson:BN-C154-petbtha}



\begin{quote}
Before the new one: say the Bengali for a cough, then the Bengali for a headache.
\end{quote}

\subsection*{You'll want to know: \bn{পেটব্যথা}}
\textbf{\bn{পেটব্যথা}} — \emph{peṭbæthā} — \textquotedblleft{}a stomach ache\textquotedblright{}.

\begin{grammarlens}[title={What you've built}]
One more word for this chapter.
\end{grammarlens}

\subsection*{Your turn}
\begin{itemize}
\raggedright
  \item \emph{Say it:} \emph{peṭbæthā}
  \item \emph{Say it:} \emph{peṭbæthā}, once more
  \item \emph{Say it:} say \emph{peṭbæthā}
  \item \emph{Recall:} say the Bengali for a cough, then the Bengali for a headache, then say \emph{peṭbæthā} again
\end{itemize}

\begin{tcolorbox}[breakable,colback=teal!4,colframe=teal!35!black,title={Before you move on}]
What does \bn{পেটব্যথা} mean? (\textquotedblleft{}A stomach ache\textquotedblright{}.) Say it once more.
\end{tcolorbox}
